% !PRINTED 327
\weberpage{327}

%<<CONT
Liegt $x$ nahe an einer Wurzel, so giebt einer der beiden Werthe
von $X$ einen noch genaueren Werth dieser Wurzel.

\begin{center}
\rule{2cm}{0.4pt}
\end{center}

Ebenso wie durch $X_1, X_2$ ein Intervall bestimmt ist, in dem
gar keine Wurzel der Gleichung liegt, so kann man auch ein
Intervall bestimmen, in dem gewiss Wurzeln liegen. Auch hierbei
geht man am besten von der homogenen Form aus und benutzt
die lineare Transformation.

Sei also wie vorher $f(x, y)$ eine Form $n^{\text{ten}}$ Grades mit lauter
reellen Wurzeln; und seien ferner $x:y$, $\xi:\eta$, $\xi':\eta'$ drei reelle
Werthe, von denen zwei willkürlich sind, und von denen der
dritte durch die Relation

\weq{16}{\begin{split}
{}[\xi f'(x) + \eta f'(y)]\,[\xi' f'(x) + \eta' f'(y)] \\
{}- H(\xi y - \eta x)\,(\xi' y - \eta' x) = 0
\end{split}}

von den beiden anderen abhängt. Die Gleichung (16) bleibt
ungeändert, wenn wir die drei Grössenpaare $x, y$; $\xi, \eta$; $\xi', \eta'$
gleichzeitig derselben linearen Transformation unterwerfen. Daher
können alle Schlüsse, die wir aus einer speciellen Form dieser
Gleichung ziehen, verallgemeinert werden. Die lineare Substitution
lässt sich so bestimmen, dass für die neuen Variablen $y = 0$,
$\xi = 0$ wird; dann ergiebt (16) für $\xi'$, $\eta'$ die Gleichung

\[
a_1(n\,a_0\,\xi' + a_1\,\eta') = [2\,n\,a_0\,a_2 - (n-1)\,a_1^2]\eta',
\]

oder also, wenn man $a_0 = 1$, $\eta' = 1$ annimmt,

\[
\xi' = -\frac{a_1^2 - 2\,a_2}{a_1}.
\]

Es ist aber, wenn wir, wie oben, mit $\alpha$ die Wurzeln von $f$
bezeichnen

\[
a_1^2 - 2\,a_2 = S(\alpha^2), \quad a_1 = -S(\alpha);
\]

also wird

\[
\xi' = \frac{S(\alpha^2)}{S(\alpha)}.
\]

Dafür lässt sich auch setzen:

\[
S[\alpha(\alpha - \xi')] = 0,
\]

und daraus ergiebt sich, da alle $\alpha$ reell sein sollen, dass die
Producte $\alpha(\alpha - \xi')$ weder alle positiv noch alle negativ sein
können, oder mit anderen Worten, es muss ein Theil der Wurzeln
%CONT>>
